\documentclass[11pt]{article}
\usepackage[margin=0.9in]{geometry}
\usepackage{amsmath,amssymb,amsthm}
\usepackage[T1]{fontenc}
\usepackage{lmodern}
\usepackage[hidelinks]{hyperref}
\usepackage{microtype}
\newtheorem{theorem}{Theorem}
\newtheorem{lemma}{Lemma}
\newcommand{\Ang}{\mathcal A}
\title{Disproof of Conjecture 00000001569\\\large Arbitrarily large point sets with at most two angles}
\author{C0ldSmi1e}
\date{9 October 2026}
\begin{document}
\maketitle
\begin{abstract}
The conjecture asserts a uniform lower bound of $n^2-O(n)$ for the number
of distinct angles determined by an $n$-point set. Neither language of the
original requires noncollinearity. We give an explicit $n$-point subset of
the Euclidean plane with at most two angle values, for every $n$. These
sets violate the claimed bound for every linear error constant and every
asymptotic threshold. Lean 4 formalizes the actual Euclidean angles and
the full quantified negation.
\end{abstract}

\section{Statement and conventions}
The English original states: ``The angle set of a point set is the set of
all angles it determines. Conjecture: An n-point set determines at least
n\textsuperscript{2} $-$ O(n) distinct angles; and the lower bound is attained
by the equidistribution of arguments of lattice families.'' The Chinese
version has the same lower-bound assertion and supplies no additional
restriction on the point sets. Both are included in \texttt{ORIGINAL.md}.

We use the ordinary planar Euclidean interpretation. For a finite set
$S\subset\mathbb R^2$, let
\[
 \Ang(S)=\{\angle abc : a,b,c\in S\text{ are pairwise distinct}\}.
\]
Here $b$ is the vertex and $\angle abc\in[0,\pi]$ is the usual unsigned
Euclidean angle between $a-b$ and $c-b$, in radians. Equal angle values
are counted only once. All rays are nonzero. We include the collinear
values $0$ and $\pi$; a convention excluding these values can only decrease
our counterexamples' angle counts. The source does not explicitly name the
ambient dimension or orientation; our convention is stated rather than
silently built into the argument. No general-position condition is assumed.

The standard uniform meaning of the lower-bound clause is
\begin{equation}\label{eq:claim}
 \exists C\in\mathbb R_{\ge0}\ \exists N\in\mathbb N\
 \forall n\ge N\ \forall S\subset\mathbb R^2\text{ finite},\quad
 |S|=n\ \Longrightarrow\ n^2-Cn\le |\Ang(S)|.
\end{equation}
The constant $C$ and threshold $N$ are independent of both $n$ and $S$.
We disprove this clause. The additional lattice-attainment phrase is not
precisely defined in the source and is not given an invented formalization:
the failure of a necessary conjunct already refutes the stated conjunction.

\section{The counterexample family}
For each natural number $n$, set
\[
 S_n=\{(j,0):j\in\mathbb Z,\ 0\le j<n\}.
\]
The first coordinate makes these $n$ points distinct, so $|S_n|=n$.
Every three points are collinear. At any vertex of a pairwise-distinct
triple, the two nonzero rays either point in the same direction or in
opposite directions. Their angle is respectively $0$ or $\pi$. Therefore
\begin{equation}\label{eq:two}
 \Ang(S_n)\subseteq\{0,\pi\},\qquad |\Ang(S_n)|\le2
 \quad\text{for every } n\in\mathbb N.
\end{equation}
This includes $n<3$, where there are no qualifying triples.

\begin{theorem}
For every real number $C$ and every $N\in\mathbb N$, there exist
$n\ge N$ and a finite planar set $S$ with $|S|=n$ such that
\[
 |\Ang(S)|<n^2-Cn.
\]
Consequently, \eqref{eq:claim} is false.
\end{theorem}
\begin{proof}
By the Archimedean property, choose $m\in\mathbb N$ with
$m>\max(C,2)$. Put $r=\max(N,m)$ and $n=r+1$. Then $n\ge N$,
$r>C$ and $r>2$, hence $n-C>1$ and $n>3$. It follows that
\[
 n^2-Cn=n(n-C)>n>3>2\ge |\Ang(S_n)|.
\]
Taking $S=S_n$ proves the assertion. In particular it provides a strict
violation for any nonnegative $C$ and any threshold $N$ proposed in
\eqref{eq:claim}, which is its negation.
\end{proof}

\section{Formalization and verification scope}
The project pins Lean 4.19.0 and Mathlib v4.19.0. The manifest records
every dependency revision. In \texttt{Solution.lean},
\texttt{Plane} is \texttt{EuclideanSpace} over the reals in dimension two.
The finite set \texttt{angleSet} takes the image of all pairwise-distinct
triples under Mathlib's \texttt{EuclideanGeometry.angle}; the theorem
\path|mem_angleSet_iff| proves its precise membership characterization.
Thus the counted values are genuine angles, not a discrete surrogate.

The declarations \path|configuration_card| and
\path|configuration_angleSet_card_le| establish the size and
angle-count bounds. The geometry step proves the x-axis collinear and
uses Mathlib's collinearity characterization by angles $0$ or $\pi$.
The theorem \path|arbitrarily_large_counterexamples| proves the full
statement above for every real $C$; \path|not_uniformQuadraticLowerBound|
negates \eqref{eq:claim}. Finally, \path|original_conjunction_false|
shows that adjoining any proposition for the attainment clause cannot
make the conjunction true. It makes no independent claim about a
particular lattice-equidistribution mechanism.

The submission supplies the full Lean source, pinned project, axiom
inspection, reproducible verification instructions, and validation
records. The proof is deductive and requires no numerical experiment or
external computation. Local verification and independent semantic review
are reported separately from any subsequent maintainer acceptance.

\paragraph{Source identity.}
The original is \texttt{conjectures/00000001569.md} at the repository commit below; the second identifier is its Git blob.
\begin{quote}\small\ttfamily
9b795e7a94a6076e49a65a6489e1caf9153abd23\\
ec72f71c9bee352ef1bcb0201554f357dc7ff480
\end{quote}
\end{document}
